% Unit 5 free-response set -- 2 long (10) + 3 short (4) = 32 pts.
\documentclass{apchem-exam}

\apcunit{5}
\apcunittitle{Kinetics}
\apcdoctitle{Free-Response Set}
\apcsubtitle{CED 5.1--5.11 \textbullet\ 5 questions \textbullet\ 32 points \textbullet\ 50 minutes}

\begin{document}
\apctitleblock

\begin{apcnote}[Scope --- one CED exclusion applies]
Covers \ced{5.1} through \ced{5.11}, worth
\SI{7}{\percent}--\SI{9}{\percent} of the exam.

Topic \ced{5.6} excludes \textbf{calculations involving the Arrhenius
equation}. You will not be asked to extract an activation energy from two
rate constants or from a plot of $\ln k$ against $1/T$. Reading an energy
profile and reasoning qualitatively about temperature and catalysts stay
fully examinable, and both appear below.
\end{apcnote}

\examsection{II}{Free Response}{5 questions \textbullet\ 32 points \textbullet\ 50 minutes}{\calculatorok}

% =====================================================================
\frq{long}{10}{\ced{5.1} \ced{5.2} \ced{5.3} \textbullet\ rate law from initial rates}

For \ce{2NO(g) + O2(g) -> 2NO2(g)}, the following initial rates were
measured at constant temperature.

\begin{center}
\begin{tabular}{cccc}
\hline
\textbf{Experiment} & $[\ce{NO}]_0$ (M) & $[\ce{O2}]_0$ (M) &
\textbf{Initial rate} (M/s) \\
\hline
1 & 0.020 & 0.010 & $2.5\times10^{-5}$ \\
2 & 0.040 & 0.010 & $1.0\times10^{-4}$ \\
3 & 0.020 & 0.020 & $5.0\times10^{-5}$ \\
\hline
\end{tabular}
\end{center}

\begin{parts}
  \item Determine the order with respect to each reactant, showing your
        reasoning. \workspace[3cm]
        \keyonly{\begin{apcanswer}\textbf{\ce{NO}:} compare 1 and 2, where
        $[\ce{O2}]$ is constant. Doubling $[\ce{NO}]$ multiplies the rate
        by $1.0\times10^{-4}/2.5\times10^{-5} = 4$. Since $2^n = 4$,
        $n = \mathbf{2}$.

        \textbf{\ce{O2}:} compare 1 and 3, where $[\ce{NO}]$ is constant.
        Doubling $[\ce{O2}]$ doubles the rate, so
        $n = \mathbf{1}$.\end{apcanswer}}
  \item Write the rate law and state the overall order.
        \answerlines[2]{$\text{rate} = k[\ce{NO}]^2[\ce{O2}]$; overall
        order $= \mathbf{3}$.}
  \item Calculate $k$, with units. \workspace[2.8cm]
        \keyonly{\begin{apcanswer}From experiment 1:

        $k = \dfrac{2.5\times10^{-5}}{(0.020)^2(0.010)} =
        \dfrac{2.5\times10^{-5}}{4.0\times10^{-6}} = \mathbf{6.3}$

        Units: rate is \si{\Molar\per\second} and
        $[\ce{NO}]^2[\ce{O2}]$ is \si{\Molar\cubed}, so $k$ is in
        \textbf{M\textsuperscript{-2}\,s\textsuperscript{-1}}.

        Check with experiment 2: $6.25(0.040)^2(0.010) =
        1.0\times10^{-4}$ \checkmark\end{apcanswer}}
  \item Explain why the order with respect to \ce{NO} cannot be assumed
        from its coefficient of 2 in the balanced equation.
        \answerlines[3]{Orders are experimental quantities set by the
        \emph{mechanism} --- in particular by the slowest step --- not by
        the overall stoichiometry. The two agree only when the reaction is
        genuinely elementary, which cannot be assumed. Here they happen to
        coincide, but that had to be established from the data.}
\end{parts}

\begin{rubric}
  \rpoint{4}{(a) 2 pts order 2 in \ce{NO} with the comparison shown; 2 pts
    order 1 in \ce{O2}. Holding the other reactant constant must be
    explicit for full credit}
  \rpoint{2}{(b) 1 pt the rate law; 1 pt overall order 3}
  \rpoint{2}{(c) 1 pt $k = \num{6.3}$; 1 pt correct units. A bare number
    earns 1}
  \rpoint{2}{(d) 1 pt orders are experimental; 1 pt they follow the
    mechanism, and coefficients may be used only for an elementary step}
\end{rubric}

% =====================================================================
\frq{long}{10}{\ced{5.4} \ced{5.5} \ced{5.7}--\ced{5.10} \ced{5.11} \textbullet\ mechanism, profile and catalysis}

A proposed mechanism for \ce{2N2O5 -> 4NO2 + O2} is
\begin{align*}
\text{Step 1:}\quad & \ce{N2O5 -> NO2 + NO3} \qquad \text{(slow)}\\
\text{Step 2:}\quad & \ce{NO2 + NO3 -> NO + O2 + NO2} \qquad \text{(fast)}\\
\text{Step 3:}\quad & \ce{NO + N2O5 -> 3NO2} \qquad\quad \text{(fast)}
\end{align*}

\begin{parts}
  \item Identify the intermediates and state how you recognized them.
        \answerlines[3]{\ce{NO3} and \ce{NO}. Each is \emph{produced} in
        one step and \emph{consumed} in a later one, so both appear in the
        mechanism but not in the overall equation.}
  \item Write the rate law this mechanism predicts, with justification.
        \answerlines[3]{The overall rate cannot exceed that of the slowest
        step, so the rate law comes from step 1. That step is
        \textbf{elementary}, so its order may be read from its own
        coefficients:

        $\text{rate} = k[\ce{N2O5}]$

        First order, even though the overall equation shows a coefficient
        of 2.}
  \item Describe how the energy profile for this three-step mechanism
        differs from that of a single-step reaction, and state which
        barrier is highest. \answerlines[3]{It shows \textbf{three}
        maxima, one per step, separated by shallow minima where the
        intermediates sit. The \textbf{first} barrier is the highest,
        because step 1 is the slow, rate-determining step.}
  \item A catalyst is introduced. State its effect on the rate, on
        $\Delta H$, and on the position of any equilibrium.
        \workspace[2.6cm]
        \keyonly{\begin{apcanswer}\textbf{Rate: increases.} The catalyst
        offers an alternative pathway with a \emph{lower activation
        energy}, so a larger fraction of collisions has enough energy to
        react.

        \textbf{$\Delta H$: unchanged.} Reactants and products sit at the
        same energies as before; only the barrier between them is lowered.

        \textbf{Equilibrium position: unchanged.} The forward and reverse
        reactions are accelerated equally, so equilibrium is reached
        sooner but lies in exactly the same place.\end{apcanswer}}
\end{parts}

\begin{rubric}
  \rpoint{2}{(a) 1 pt both intermediates; 1 pt produced-then-consumed}
  \rpoint{3}{(b) 1 pt rate law taken from the slow step; 1 pt
    $k[\ce{N2O5}]$; 1 pt noting coefficients are usable only because the
    step is elementary}
  \rpoint{3}{(c) 1 pt three maxima with intermediates between; 1 pt the
    first is highest; 1 pt linking that to the slow step}
  \rpoint{2}{(d) 1 pt lower activation energy via an alternative path;
    1 pt $\Delta H$ and equilibrium position both unchanged. ``It lowers
    the energy of the reactants'' earns 0}
\end{rubric}

% =====================================================================
\frq{short}{4}{\ced{5.3} \textbullet\ first-order half-life}

A first-order decomposition has $k = \SI{1.5e-3}{\per\second}$ and an
initial concentration of \SI{0.80}{\Molar}.

\begin{parts}
  \item Calculate the half-life. \answerlines[2]{$t_{1/2} =
        \dfrac{0.693}{1.5\times10^{-3}} = \mathbf{4.6\times10^{2}}$~s}
  \item State the concentration remaining after three half-lives, and
        explain why each successive half-life takes the same time.
        \answerlines[3]{$0.80 \to 0.40 \to 0.20 \to \mathbf{0.10}$~M.

        For a first-order reaction the half-life depends only on $k$, not
        on concentration: as the concentration falls the rate falls in
        exact proportion, so every halving takes the same interval.}
\end{parts}

\begin{rubric}
  \rpoint{2}{(a) 1 pt the relationship; 1 pt \SI{4.6e2}{\second}}
  \rpoint{2}{(b) 1 pt \SI{0.10}{\Molar}; 1 pt half-life independent of
    concentration for first order}
\end{rubric}

% =====================================================================
\frq{short}{4}{\ced{5.1} \textbullet\ relating rates of different species}

In \ce{2NO(g) + O2(g) -> 2NO2(g)}, oxygen is consumed at
\SI{0.0075}{\Molar\per\second}.

\begin{parts}
  \item Determine the rate at which \ce{NO} is consumed and \ce{NO2}
        formed. \answerlines[2]{Rates scale with the coefficients:
        \ce{NO} is consumed at $2(0.0075) = \mathbf{0.015}$~M/s and
        \ce{NO2} formed at $2(0.0075) = \mathbf{0.015}$~M/s.}
  \item Explain why dividing each concentration change by its coefficient
        gives the same number for every species.
        \answerlines[2]{Because the coefficients fix the fixed proportions
        in which the species are consumed and produced. Dividing out those
        proportions leaves a single rate describing the reaction as a
        whole, independent of which species was monitored.}
\end{parts}

\begin{rubric}
  \rpoint{2}{(a) 1 pt each rate, with signs or wording showing consumption
    versus formation}
  \rpoint{2}{(b) 1 pt dividing by coefficients; 1 pt gives one rate for
    the reaction}
\end{rubric}

% =====================================================================
\frq{short}{4}{\ced{5.5} \ced{5.6} \textbullet\ the collision model}

\begin{parts}
  \item State the two conditions a collision must meet to produce
        reaction. \answerlines[2]{Energy at least equal to the activation
        energy, and correct \textbf{orientation} of the colliding
        molecules.}
  \item A rise of \SI{10}{\celsius} can roughly double a reaction rate,
        although it raises the average molecular speed by only a few
        percent. Explain. \workspace[2.8cm]
        \keyonly{\begin{apcanswer}What matters is the \emph{fraction} of
        collisions carrying at least the activation energy, and that
        fraction lies far out in the high-energy tail of the distribution
        of molecular energies.

        Warming shifts and broadens the whole distribution. Near the
        average that is a small change, but in the tail the population
        rises steeply, so the number of molecules above the threshold can
        easily double. The increase in collision \emph{frequency} is minor
        by comparison.\end{apcanswer}}
\end{parts}

\begin{rubric}
  \rpoint{2}{(a) 1 pt each condition}
  \rpoint{2}{(b) 1 pt the fraction above the activation energy rises
    sharply; 1 pt this outweighs the change in collision frequency.
    ``Molecules move faster and collide more'' alone earns 1. Per the
    \ced{5.6} exclusion, no Arrhenius calculation may be required}
\end{rubric}

\examtotal

\end{document}
